\documentclass{exam}
\usepackage{amsmath, amssymb}

\pagestyle{headandfoot}
\firstpageheadrule
\runningheadrule
\firstpageheader{Prof. Rong \\ Introduction to Real Analysis I}{Homework 5}{Aadhithya Saravanan}
\runningheader{Introduction to Real Analysis I \\ Homework 5}{}{Saravanan}
\firstpagefooter{}{}{}
\runningfooter{ }{\thepage}{ }

\printanswers

\begin{document}

\underline{Exercise 2.7.1}
\newline
\begin{parts}
    \part By showing that $(s_n)$ is Cauchy, we will know that it is convergent. For $(s_n)$ to be Cauchy, there must exist, for every $\epsilon>0$, an $N\in\mathbb{N}$ such that for any $n>m\ge N$, it follows that $|s_n-s_m|=|a_{m+1}-a_{m+2}+a_{m+3}-a_{m+4}+\cdots\pm a_n|<\epsilon$. Since $(a_n)\to 0$ and the sequence is monotone decreasing, every term in $(a_n)$ must be at least 0. By grouping the terms in $|a_{m+1}-a_{m+2}+a_{m+3}-a_{m+4}+\cdots\pm a_n|$, we have $|(a_{m+1}-a_{m+2})+(a_{m+3}-a_{m+4})+\cdots|$. Because the sequence is monotone decreasing, we have $|(a_{m+1}-a_{m+2})+(a_{m+3}-a_{m+4})+\cdots|\ge|(0)+(0)+\cdots|=0$, so $|s_n-s_m|\ge0$. By grouping in a different way, we have $|a_{m+1}-a_{m+2}+a_{m+3}-a_{m+4}+\cdots\pm a_n|=|a_{m+1}-(a_{m+2}-a_{m+3})-\cdots|\le|a_{m+1}-(0)-(0)-\cdots|\le a_{m+1}$. So, $0\le|s_n-s_m|\le a_{m+1}$. Since $(a_n)\to 0$, for every $\epsilon>0$, there exists an $L\in\mathbb{N}$ such that $l\ge L$ implies $|a_l|<\epsilon$. We can apply this by choosing $N=L$. We also know $m+1>m\ge L$, so we have $0\le|s_n-s_m|\le a_{m+1}<\epsilon$. So, the Cauchy Criterion is satisfied and the Alternating Series Test is proven.
    \part The even partial sums, which we can call $(s_{2n})$, are increasing, while the odd partial sums, which we can call $(s_{2n+1})$, are decreasing. This can be shown by induction. The first and second even partial sums are $s_{2}=a_1-a_2$ and $s_{4}=a_1-a_2+a_3-a_4$. $s_{4}-s_{2}=(a_1-a_2+a_3-a_4)-(a_1-a_2)=a_3-a_4\ge0$. Assuming the $k$th even partial is greater than or equal to the $k-1$th even partial sum, we need to show that $s_{2(k+1)}-s_{2k}\ge0$. Since the difference between the partial sums are just the terms $a_{2k-1}-a_{2k}$, which is at least 0 by monotonicity, the statement holds and the even partial sums are increasing. The same method with a different grouping method can be used to show that the odd partial sums are decreasing. We also know that adding one extra term to an even partial sum adds a positive number, increasing it. So, each even partial sum is less than or equal to the next odd partial sum. Adding the next term, which is negative, to an odd partial sum decreases the sum, so each odd partial sum is greater than or equal to the next even partial sum. Let $I_n=[s_{2n},s_{2n+1}]$. Since the even partial sums are increasing and the odd partial sums are decreasing, the lower bound on each interval is increasing while the upper bound is decreasing. Therefore, $I_{n+1}\subseteq I_n$. Because the sequence $(a_n)\to0$, the length of the intervals decrease to 0. By the Nested Interval Property, there must exist an element in the intersection of all $I_n$. Since the length of the intervals go to 0, this element is unique. So, the partial sums converge to this element, proving the series converges.
    \part The elements of the subsequence $(s_{2n})$ are the even partial sums of the series, which we have shown is an increasing sequence in the last part. The first partial sum, which is just the $a_1$ term, is an upper bound for the even partial sums. Each even partial sum can be written as $s_{2n}=a_1-(a_2-a_3)-(a_4-a_5)-\cdots-(a_{2n-2}-a_{2n-1})-a_{2n}\le a_1-(0)-(0)-\cdots-(0)-a_{2n}\le a_1$. Since $(s_{2n})$ is increasing and bounded, by the Monotone Convergence Theorem, it converges. The same can be shown for $(s_{2n+1})$ with a similar method, so it is also convergent. Let $L_1$ be the limit of the even partial sums and $L_2$ be the limit of the odd partial sums. Since $s_{2n+1}=s_{2n}+a_{2n+1}$, we can take the limits of this equation as $n\to\infty$. We get $L_2=L_1+0$, so the limits are equal. Therefore, the entire series converges to the same limit.
    \newline
\end{parts}

\underline{Exercise 2.7.2}
\newline
\begin{parts}
    \part For all $n\in\mathbb{N}$, $0\le\frac{1}{2^n+n}\le\frac{1}{2^n}$, so we can use the Comparison Test. $\sum_{n=1}^{\infty}\frac{1}{2^n}=\sum_{n=0}^{\infty}(\frac{1}{2})^n-(\frac{1}{2})^0=\frac{1}{1-\frac{1}{2}}-1=1$. Since $\sum_{n=1}^{\infty}\frac{1}{2^n}$ is convergent, $\sum_{n=1}^{\infty}\frac{1}{2^n+n}$ is convergent by the Comparison Test.
    \part If we can show that $\sum_{n=1}^{\infty}|\frac{sin(n)}{n^2}|$ is convergent, then we will have that $\sum_{n=1}^{\infty}\frac{\sin(n)}{n^2}$ is convergent as well by the Absolute Convergence Test. $|\frac{sin(n)}{n^2}|=\frac{|sin(n)|}{n^2}<\frac{1}{n^2}<\frac{1}{n(n-1)}$ for $n>2$, so we can use the Comparison Test. Let $s_N$ represent the partial sum $\sum_{n=2}^{N}\frac{1}{n(n-1)}$. This is equal to $\sum_{n=2}^{N}\frac{1}{n-1}-\frac{1}{n}=(1-\frac{1}{2})+(\frac{1}{2}-\frac{1}{3})+\dots+(\frac{1}{N-2}-\frac{1}{N-1})+(\frac{1}{N-1}-\frac{1}{N})=1-\frac{1}{N}$. As $N\to\infty$, the limit of the sequence of partial sums goes to 1, so $\sum_{n=2}^{\infty}\frac{1}{n(n-1)}$. Then, $\sum_{n=2}^{\infty}\frac{1}{n^2}$ is convergent, and adding the term for $n=1$ keeps the series convergent. Then, by the Comparison and Absolute Convergence Tests, $\sum_{n=1}^{\infty}\frac{\sin(n)}{n^2}$ is convergent.
    \part The given series is equal to $\sum_{n=1}^{\infty}\frac{(-1)^{n+1}(n+1)}{2n}$. This is equal to $\sum_{n=1}^{\infty}\frac{(-1)^{n+1}n}{2n}+\sum_{n=1}^{\infty}\frac{(-1)^{n+1}}{2n}=\frac{1}{2}\sum_{n=1}^{\infty}(-1)^{n+1}+\frac{1}{2}\sum_{n=1}^{\infty}\frac{(-1)^{n+1}}{n}$ by the Algebraic Limit Theorem for Series. For the first sum, the sequence representing the partial sum up to the $N$th term is given by $s_N=\sum_{n=1}^{N}(-1)^{n+1}=\frac{1}{2}-\frac{1}{2}(-1)^N$. This sequence oscillates and is therefore divergent, so the sum does not converge. Therefore, the entire series does not converge.
    \part The given series can be written as $\sum_{n=1}^{\infty}\frac{1}{3n-2}+\frac{1}{3n-1}-\frac{1}{3n}$. After combining terms, we have $\sum_{n=1}^{\infty}\frac{9n^2-2}{(3n-2)(3n-1)(3n)}$. Then, $\sum_{n=1}^{\infty}\frac{9n^2-2}{(3n-2)(3n-1)(3n)}\ge\sum_{n=1}^{\infty}\frac{9n^2-2}{(3n)^3}\ge\sum_{n=1}^{\infty}\frac{8n^2}{27n^3}=\frac{8}{27}\sum_{n=1}^{\infty}\frac{1}{n}$, assuming using the terms for $n\ge2$. Since the harmonic series is divergent, the original series must be divergent by the Comparison Test.
    \part We can rewrite the given series as the sum of the positive terms minus the sum of the negative terms. Therefore, the series is equal to $\sum_{n=1}^{\infty}\frac{1}{2n-1}-\sum_{n=1}^{\infty}\frac{1}{(2n)^2}$. The first sum is greater than $\sum_{n=1}^{\infty}\frac{1}{2n}$, which is equal to $\frac{1}{2}\sum_{n=1}^{\infty}\frac{1}{n}$. Since the harmonic series is divergent, a scalar multiple of it is also divergent, so $\sum_{n=1}^{\infty}\frac{1}{2n-1}$ is divergent. To prove that the sum is divergent, the second sum must be convergent so that the difference of sums is divergent. $\sum_{n=1}^{\infty}\frac{1}{(2n)^2}=\frac{1}{4}\sum_{n=1}^{\infty}\frac{1}{n^2}$. This sum has been shown to converge in part (b) with a comparison to $\sum_{n=1}^{\infty}\frac{1}{n(n-1)}$. Therefore, the scalar multiple of the sum is convergent, so the difference of the sums is divergent. Therefore, the series is divergent.
    \newline
\end{parts}

\underline{Exercise 2.7.4}
\newline
\begin{parts}
    \part Let $x_n=(-1)^{n+1}$. The partial sum of the terms in sequence oscillate between 1 and 0, so it is divergent. Let $y_n=\frac{1}{n}$. This is the harmonic series, which is known to diverge. When multiplying each $x_n$ by the corresponding $y_n$, we have the sequence $x_ny_n=\frac{(-1)^{n+1}}{n}$. The sum of this sequence is the alternating harmonic series, which is convergent.
    \part Let $x_n=\frac{(-1)^{n+1}}{n}$. The sum of this sequence is the alternating harmonic series, which is convergent. Let $y_n=(-1)^{n+1}$. An upper bound for $(y_n)$ is 1, as the only distinct terms in the sequence are $-1$ and 1. When multiplying each $x_n$ by the corresponding $y_n$, we have the sequence $x_ny_n=\frac{(-1)^{n+1}(-1)^{n+1}}{n}=\frac{(-1)^{2n+2}}{n}=\frac{1}{n}$. The sum of this sequence is the harmonic series, which is divergent.
    \part This is impossible. $\sum(x_n+y_n)=\sum x_n+\sum y_n$. Let $\sum(x_n+y_n)=C$ for some $C\in\mathbb{R}$, since this sum is convergent. Let $\sum(x_n)=A$ for some $A\in\mathbb{R}$, since this sum also is convergent. Then, $\sum y_n=\sum(x_n+y_n)-\sum x_n=C-A$. Since $C-A\in\mathbb{R}$, $\sum y_n$ must be convergent. Therefore, $\sum y_n$ cannot be divergent.
    \part Let $x_n=\frac{1}{n}$ for even $n$ and $x_n=0$ for odd $n$. Then, the alternating series of this sequence is $0+\frac{1}{2}+0+\frac{1}{4}+0+\frac{1}{6}+\cdots=\frac{1}{2}(1+\frac{1}{2}+\frac{1}{3}+\cdots)$. Since this series is half of the harmonic series, which is divergent, it must also be divergent.
    \newline
\end{parts}

\underline{Exercise 2.7.7}
\newline
\begin{parts}
    \part By the definition of limits, $\lim(na_n)=l$ implies that for every $\epsilon>0$, there exists an $N\in\mathbb{N}$ such that $n\ge N$ implies $|na_n-l|<\epsilon$. This can be written as $-\epsilon<na_n-l$, which is equivalent to $\frac{-\epsilon+l}{n}<a_n$. Now, we choose $\epsilon=\frac{l}{2}$. We know $l>0$ from $\lim(na_n)=l$ because $n>0$ and $a_n>0$. Then, $\frac{l}{2n}=\frac{l}{2}\cdot\frac{1}{n}<a_n$. Since $l>0$, we know $\frac{l}{2}>0$. Then, we can compare $a_n$ to the sequence that contains each term in the harmonic series multiplied by $\frac{l}{2}$. Since that sequence is each term in the harmonic series, which is divergent, multiplied by a positive number, it must also diverge. Then, by the Comparison Test, $\sum a_n$ diverges.
    \part Let $\lim(n^2a_n)=l$. By the definition of limits, $\lim(n^2a_n)=l$ implies that for every $\epsilon>0$, there exists an $N\in\mathbb{N}$ such that $n\ge N$ implies $|n^2a_n-l|<\epsilon$. This can be written as $n^2a_n-l<\epsilon$, which is equivalent to $a_n<\frac{\epsilon+l}{n^2}$. Choose $\epsilon=1$. Then, $a_n<\frac{l+1}{n^2}=2l\cdot\frac{1}{n^2}$. This is equal to a positive constant, $l+1$, multiplied by the sequence of the reciprocals of the squares of the natural numbers, which is convergent by Corollary 2.4.7. Since multiplying by a constant does not change whether a series diverges or converges, the sum of the sequence converges. Since $a_n$ is less than each term in the sequence and is bounded below by 0, it must also converge by the Comparison Test.
    \newline
\end{parts}

\underline{Exercise 2.7.11}
\newline
\newline
We want $(a_n)$ and $(b_n)$ to take turns being the term that is equal to $\min\{a_n,b_n\}$ for each $n\in\mathbb{N}$. At any given $n$, we want $\min\{a_n,b_n\}=\frac{1}{n^2}$. At any given index, either $a_n$ or $b_n$ should equal $\frac{1}{n^2}$, while the other has a greater value. We will start with the first "block". First, let $a_1=\frac{1}{1^2}=1$ and $b_1=1$. Because the sum of just the first term in $(b_n)$ is at least 1, the sequences swap roles for $n>1$ and we move to the next block. For block 2, let $b_n=\frac{1}{n^2}$ and $a_n=a_{n-1}$ until the sum of the first terms in $a_n$ add up to at least 1. In this case, we have $a_2=1$, so the second block also has length 1. $b_2=\frac{1}{2^2}=\frac{1}{4}$. Then we again swap the roles of the sequences and move to the next block. So, $a_n=\frac{1}{n^2}$ and $b_n=b_{n-1}=\frac{1}{4}$ for $n>2$. Now, it takes 4 terms of $b_n$ to add up to at least 1, at which point $n$ is at 7, so $a_7=\frac{1}{7^2}$. We can continue this process infinitely, as each block has one of the sequences at a constant value that is positive, which means a finite number of them adds up to 1 eventually, at which point the next block starts. $(a_n)$ and $(b_n)$ are decreasing. The sequences are either strictly decreasing when they are defined as $\frac{1}{n^2}$ or are equal to their previous terms. They are also positive for any given $n$. For odd-numbered blocks, $a_n=\frac{1}{n^2}$ and $b_n=b_{n-1}$. For even-numbered blocks, $a_n=a_{n-1}$ and $b_n=\frac{1}{n^2}$. So, $\min\{a_n,b_n\}=\frac{1}{n^2}$ for $n\in\mathbb{N}$. Since this series is convergent by Corollary 2.4.7, $\sum\min\{a_n,b_n\}$ is convergent.

\end{document}